\documentclass[11pt,a4paper]{coverletter}

\senderblock
  {Radheem Bin Razi}
  {Distributed Systems \& Cloud-Native Engineer}
  {Ilmenau, Germany \textbar\ \href{mailto:sheikh.radheem@gmail.com}{sheikh.radheem@gmail.com} \textbar\ \href{https://radheem.github.io/my-cv}{portfolio}}

\begin{document}

%% ===== ENGLISH =====
\recipient{Sievert SE}{Hiring Team}{}

\opening{Dear Hiring Team,}

\vspace{8pt}\noindent\textbf{1. Warum Sievert SE?}\par\vspace{3pt}

Saubere Datenflüsse sind das unsichtbare Rückgrat einer zuverlässigen Lieferkette. Genau diesen stillen Aufbau von Operations-Transparenz finde ich an der Schnittstelle von Logistik und Produktion am spannendsten. Mit der Stelle als Data Engineer bei Sievert SE möchte ich diese operative Transparenz in Ihrem Umfeld für nachhaltige Baustoffe und intelligente Logistik verankern.

\vspace{8pt}\noindent\textbf{2. Warum ich?}\par\vspace{3pt}

In meinen bisherigen Projekten habe ich mich darauf spezialisiert, komplexe Datenströme in stabile, wartbare Pipelines zu überführen. Bei einem umfangreichen Testfeld habe ich eine Kafka-basierte Datenpipeline aufgebaut, die Echtzeitmetriken zuverlässig in InfluxDB und MongoDB überführte. Dabei lag der Fokus stets auf strenger Datenqualität und konsistenter Modellierung. So konnten nachgelagerte Analysen auf verlässlichen Grundlagen aufbauen. Ich habe diese Systeme regelmäßig gewartet, technische Prozesse dokumentiert und eng mit Anwendern zusammengearbeitet. Gemeinsam optimierten wir die Abfrageperformance und sicherten die langfristige Datenkonsistenz. Ihre Anforderungen an Reporting \& Analyse und Data Governance entsprechen genau meinem Arbeitsstil.

Ein weiteres Projekt bestand aus einer dokumentenbasierten Ingestionspipeline, die mit automatischen Wiederholungslogiken auch bei Systemausfällen keine Daten verliert. Ich pflege diese Architekturen kontinuierlich weiter und bringe fundierte Kenntnisse in Python, Cloud-Infrastrukturen und Datenmodellierung mit. Genau diese Erfahrung mit zuverlässigen ETL-Prozessen und klarer Datenarchitektur möchte ich in Ihren Aufbau des Produktionsmanagementsystems einbringen.

\vspace{8pt}\noindent\textbf{3. Warum jetzt?}\par\vspace{3pt}

Derzeit suche ich genau diese Phase, in der ich meine Erfahrung mit verteilten Datenplattformen in ein wachsendes, mittelständisches Unternehmen einbringen kann. Ich bin ab sofort vollzeitlich verfügbar. Ich kann innerhalb von zwei bis drei Wochen an jeden Standort in Deutschland wechseln. Als Inhaber einer Blue Card benötige ich keine Arbeitsgenehmigung. Der Einstieg ist daher ohne bürokratische Verzögerung möglich.

\closing{Sincerely,}{Radheem Bin Razi}

\clearpage
\selectlanguage{ngerman}

%% ===== DEUTSCH =====
\recipient{Sievert SE}{Personalabteilung}{}

\opening{Sehr geehrtes Hiring-Team,}

\vspace{8pt}\noindent\textbf{1. Warum Sievert SE?}\par\vspace{3pt}

Saubere Datenflüsse bilden das unsichtbare Rückgrat einer zuverlässigen Lieferkette. Genau diesen stillen Aufbau von Operations-Transparenz finde ich an der Schnittstelle von Logistik und Produktion am spannendsten. Mit der Position als Data Engineer bei Sievert SE möchte ich diese operative Transparenz in Ihrem Umfeld für nachhaltige Baustoffe und intelligente Logistik verankern.

\vspace{8pt}\noindent\textbf{2. Warum ich?}\par\vspace{3pt}

In meinen bisherigen Projekten habe ich mich darauf spezialisiert, komplexe Datenströme in stabile, wartbare Pipelines zu überführen. Für ein umfangreiches Testfeld habe ich eine Kafka-basierte Datenpipeline aufgebaut, die Echtzeitmetriken zuverlässig in InfluxDB und MongoDB überführte. Dabei lag der Fokus stets auf strenger Datenqualität und konsistenter Modellierung. So konnten nachgelagerte Analysen auf verlässlichen Grundlagen aufbauen. Ich habe diese Systeme regelmäßig gewartet, technische Prozesse dokumentiert und eng mit den Anwendern zusammengearbeitet. Gemeinsam optimierten wir die Abfrageperformance und sicherten die langfristige Datenkonsistenz. Ihre Anforderungen an Reporting \& Analyse und Data Governance entsprechen genau meinem Arbeitsstil.

Ein weiteres Projekt bestand aus einer dokumentenbasierten Ingestionspipeline, die dank automatischer Wiederholungslogiken auch bei Systemausfällen keine Daten verliert. Ich pflege derartige Architekturen kontinuierlich weiter und bringe fundierte Kenntnisse in Python, Cloud-Infrastrukturen und Datenmodellierung mit. Genau diese Erfahrung mit zuverlässigen ETL-Prozessen und einer klaren Datenarchitektur möchte ich in Ihren Aufbau des Produktionsmanagementsystems einbringen.

\vspace{8pt}\noindent\textbf{3. Warum jetzt?}\par\vspace{3pt}

Derzeit suche ich genau diese Phase, in der ich meine Erfahrung mit verteilten Datenplattformen in ein wachsendes, mittelständisches Unternehmen einbringen kann. Ich bin ab sofort vollzeitlich verfügbar. Ich kann innerhalb von zwei bis drei Wochen an jeden Standort in Deutschland wechseln. Als Inhaber einer Blue Card benötige ich keine zusätzliche Arbeitsgenehmigung. Der Einstieg ist daher ohne bürokratische Verzögerung möglich.

\closing{Mit freundlichen Grüßen,}{Radheem Bin Razi}

\end{document}
